\chapter{Introducción}

\section{Contexto}

En criptografía simétrica hay dos grandes familias de cifrados: los cifrados en bloque y
los cifrados de flujo. Los cifrados en bloque, como AES, dividen el mensaje en trozos de
tamaño fijo y los procesan de forma independiente. Los cifrados de flujo combinan cada bit
del mensaje con un bit de una secuencia pseudoaleatoria llamada \textit{keystream}, lo que
permite cifrar datos de cualquier longitud sin relleno y con una latencia mínima.

La idea de cifrar combinando el mensaje con una secuencia de clave bit a bit viene
de lejos. En 1917, Gilbert Vernam patentó un sistema telegráfico que hacía exactamente
eso, y en 1949 Shannon demostró que si la secuencia de clave es verdaderamente
aleatoria, tan larga como el mensaje y se usa una sola vez, el cifrado resulta
teóricamente irrompible. Es lo que se conoce como \textit{one-time pad}. El problema
es que distribuir una clave tan larga como el mensaje que se quiere proteger es
inviable en la mayoría de escenarios, así que los cifrados de flujo parten de una
clave corta (de 80 o 128 bits) y la expanden con un generador pseudoaleatorio
hasta obtener un \textit{keystream} de la longitud necesaria. La seguridad ya no es
perfecta, pero si el generador está bien diseñado, distinguir el keystream de una
secuencia verdaderamente aleatoria debería ser computacionalmente imposible.

En la práctica, la mayoría de cifrados de flujo diseñados para hardware se han basado
en LFSR. La razón es sencilla: un LFSR se implementa con flip-flops y puertas XOR,
que son los componentes más baratos que existen en un circuito digital. La teoría
matemática detrás de los LFSR está muy desarrollada (polinomios primitivos,
secuencias de longitud máxima, autocorrelación), lo que permite conocer de antemano
las propiedades de la secuencia base. El reto siempre ha sido el mismo: cómo combinar
uno o varios LFSR para que la salida final no sea lineal sin perder esas buenas
propiedades.

Los cifrados de flujo llevan décadas usándose en sistemas donde la velocidad o el coste
hardware son críticos. Los dos ejemplos más conocidos son A5/1 en la telefonía GSM y E0 en
Bluetooth. A5/1 utiliza tres LFSR de grados 19, 22 y 23, combinados con un mecanismo
de reloj irregular, que en cada ciclo, un bit de mayoría decide cuáles de los tres
registros avanzan y cuáles se quedan parados. E0 es algo distinto, combina cuatro
LFSR con una máquina de estados finita de 4 bits que mezcla las salidas. En ambos
casos la razón de elegir un cifrado de flujo frente a uno de bloque fue la misma:
generar un bit de \textit{keystream} requiere poco más que operaciones XOR, así que
los circuitos son pequeños y el consumo bajo.

Ninguno de estos dos diseños ha resistido bien el paso del tiempo. A5/1 se diseñó
en 1987 como parte del estándar GSM, pero su especificación se mantuvo en secreto
durante más de una década bajo la premisa de que la oscuridad añadiría seguridad.
En 1999, Briceno, Goldberg y Wagner consiguieron hacer ingeniería inversa del
algoritmo a partir de un teléfono GSM, y ese mismo año Biryukov, Shamir y Wagner
publicaron ataques que permitían recuperar la clave con un esfuerzo del orden de
$2^{40}$ operaciones, muy por debajo de los $2^{64}$ que se suponía que
ofrecía~\cite{biryukov2000}. Finalmente en 2010, Karsten Nohl
demostró que con tablas arcoíris precomputadas se podían descifrar llamadas GSM en
tiempo casi real usando un portátil y una radio definida por software.

E0 tampoco aguantó demasiado. En 2004, Lu y Vaudenay publicaron un ataque que
explota correlaciones condicionales entre la salida y los estados internos de los
LFSR, reduciendo la complejidad real a unas $2^{38}$ operaciones con los primeros
$2^{24}$ bits del keystream, muy lejos de los $2^{128}$ que la longitud de clave
prometía~\cite{lu2004}. El problema es que la máquina de estados que combina los
cuatro registros no introduce suficiente no linealidad, ya que hay correlaciones detectables
entre la salida y los registros individuales.

En los dos casos la causa del fallo es la misma: la construcción que combina los
LFSR no consigue romper la linealidad lo bastante, y eso es explotable.

El componente que aparece en prácticamente todos estos diseños es el \textbf{registro de
desplazamiento con retroalimentación lineal}, conocido como LFSR (\textit{Linear Feedback
Shift Register}). Un LFSR es un registro de $n$ celdas binarias cuyo siguiente estado se
calcula como una combinación lineal (XOR) de algunas posiciones del estado actual,
definidas por un \textbf{polinomio de retroalimentación}. Si el polinomio es primitivo
sobre $\mathbb{F}_2$, el registro recorre los $2^n - 1$ estados no nulos antes de
repetirse, generando secuencias con buenas propiedades estadísticas: distribución
equilibrada de ceros y unos, buena autocorrelación y periodo largo~\cite{menezes1996}.

El problema es que la linealidad del LFSR también es su mayor debilidad. El algoritmo
de Berlekamp-Massey~\cite{massey1969} permite reconstruir el polinomio de
retroalimentación y el estado completo de un LFSR de grado $n$ a partir de solo $2n$
bits de salida. Para un LFSR de grado 20, bastan los primeros 40 bits observados para romperlo por
completo, da igual que su periodo supere el millón de bits. Un LFSR solo
es inviable como generador criptográfico.

Pero Berlekamp-Massey no es la única amenaza. Cuando se combinan varios LFSR, aparecen
los \textbf{ataques de correlación}: si la salida del generador tiene dependencia
estadística con alguno de los LFSR internos, un atacante puede aislarlo y atacarlo por
separado~\cite{meier1994}. Este tipo de ataque es exactamente el que rompe el generador
de Geffe, como se verá en el capítulo~\ref{cap:generadores}.

Para resistir estos ataques se han ido proponiendo en esta memoria distintas construcciones que combinan
varios LFSR o aplican transformaciones no lineales sobre sus salidas. Una de las
primeras fue la de Geffe~\cite{geffe1973} en 1973, que usa un multiplexor controlado
por un tercer LFSR para seleccionar la salida de uno de los otros dos. Con el tiempo
se vio que era vulnerable a ataques de correlación, y se buscaron alternativas más
robustas: el generador de paso alternado o ASG~\cite{gunther1987}, el generador de
contracción~\cite{coppersmith1993} y el de autocontracción~\cite{meier1994}. En este
trabajo se estudian estas cuatro construcciones.

\section{Estado del arte}

Tras los problemas de seguridad de A5/1 y E0, quedó claro que hacía falta encontrar cifrados de flujo
nuevos y bien analizados. En 2004, la red europea de investigación ECRYPT lanzó el proyecto eSTREAM~\cite{estream2008},
un concurso abierto para seleccionar nuevos cifrados de flujo divididos en dos perfiles:
software (para procesadores de propósito general) y hardware (para circuitos con
restricciones de área y consumo). El concurso se cerró en 2008 con un portafolio de
siete algoritmos recomendados.

Lo interesante para este trabajo es que varios de los finalistas del perfil hardware
siguen usando LFSR como componente base. Grain, por ejemplo, combina un LFSR de
grado 80 con un registro de desplazamiento con retroalimentación no lineal (NLFSR) del
mismo tamaño y una función booleana de filtrado. Trivium usa tres registros de
desplazamiento de 93, 84 y 111 bits con realimentación cruzada no lineal entre ellos.
En ambos casos se mantiene el LFSR como fuente de buenas propiedades estadísticas y
de periodo, pero la seguridad la aporta la capa no lineal que se le añade encima.

Esto muestra que los LFSR no se han abandonado como componente criptográfico, sino que
la lección aprendida de A5/1 y E0 ha sido que la función que combina las salidas de los
registros es lo que determina la seguridad del conjunto. Un LFSR bien elegido garantiza
periodo largo y buenas propiedades estadísticas de base, pero si la combinación no rompe
la linealidad de forma efectiva, el generador es vulnerable.

Los generadores que se estudian en este trabajo (Geffe, contracción, autocontracción
y ASG) representan una generación anterior a eSTREAM y ninguno se usa hoy en día en
aplicaciones reales. Pero son precisamente los que aparecen en los libros de texto de
criptografía~\cite{menezes1996} como ejemplos de distintas estrategias para combinar
LFSR, y estudiarlos permite entender qué principios de diseño funcionan y cuáles no
antes de abordar construcciones más sofisticadas.

Desde el punto de vista experimental, la mayoría de publicaciones analizan cada
generador de forma aislada en el contexto de su propuesta original. Comparaciones
directas entre varios de ellos usando los mismos parámetros, midiendo complejidad
lineal y pasando los mismos tests estadísticos NIST, no son tan habituales. Este
trabajo intenta aportar esa comparación para estas cuatro construcciones concretas.

\section{Motivación}

La teoría de los LFSR está bien cubierta en la literatura, pero la mayoría de
publicaciones se centran en un único generador y lo analizan de forma aislada. Cuesta
encontrar trabajos que pongan varios generadores lado a lado, con los mismos parámetros
y los mismos tests, para comparar directamente cómo se comportan. Esa comparación es
lo que me interesaba hacer: implementar cuatro generadores clásicos, medir su
complejidad lineal con Berlekamp-Massey y pasar las secuencias por los tests del
estándar NIST SP 800-22~\cite{nist2010} para ver si las diferencias teóricas se
reflejan de verdad en la práctica.

La elección de los cuatro generadores no es arbitraria. Tres de ellos (contracción,
autocontracción y ASG) se consideran razonablemente robustos: tienen complejidad
lineal alta y no se conocen ataques de correlación prácticos contra ellos. El cuarto,
el generador de Geffe~\cite{geffe1973}, es un caso deliberadamente débil. Geffe
aparece en todos los libros de texto porque fue una de las primeras propuestas para
combinar LFSR de forma no lineal, pero tiene un problema conocido: su función de
combinación tiene correlación $0{,}75$ con dos de los tres LFSR internos, lo que
permite atacarlos por separado en vez de tener que romper el sistema completo. Incluir
un generador que se sabe que falla es lo que permite ver qué distingue a un diseño
seguro de uno vulnerable cuando se miran los resultados de los tests.

Los cuatro generadores representan estrategias muy distintas para introducir no
linealidad. Geffe usa una función booleana que selecciona entre dos LFSR según el
valor de un tercero. El generador de contracción descarta bits de un LFSR según la
salida de otro, de forma que la secuencia resultante tiene longitud variable e
impredecible. El de autocontracción lleva esta idea un paso más allá: un solo LFSR
se utiliza a la vez como fuente de bits y como mecanismo de selección. Y el ASG
alterna entre dos LFSR controlados por un tercero, produciendo un bit de uno u otro
en cada ciclo según el valor del registro de control. Comparar estas cuatro
estrategias con las mismas herramientas permite ver cuál resiste mejor y, sobre todo,
entender por qué.

\subsection*{Implementación desde cero}

Decidí implementar \textbf{todo el software desde cero en Python}, sin usar ninguna
librería externa de criptografía. Esto incluye el LFSR, el algoritmo de
Berlekamp-Massey, los cuatro generadores y los ocho tests estadísticos del estándar
NIST SP 800-22.

Sé que existen librerías como \texttt{pylfsr} para los registros de desplazamiento o
\texttt{nistrng} para los tests estadísticos. Usarlas habría simplificado bastante el
desarrollo, pero el hecho de implementar todo desde cero hace que sea mucho mas fácil
la legibilidad del código y la comprensión de lo que realmente hace cada parte.
Si se usaran librerías externas, el código de esta memoria se limitaría a llamar a
funciones de esas librerías sin que quede claro qué hay detrás de cada una.
Por ejemplo, usar una función ya implementada de Berlekamp-Massey habría dado el
resultado correcto, pero habría convertido gran parte del trabajo en una caja negra de
la que se desconocería su funcionamiento interno.

Con los tests NIST pasa algo parecido, codificar cada test te obliga a
tener claro qué hipótesis estadística se está evaluando y qué significa realmente
que una secuencia lo supere o lo falle.

Además, al haber desarrollado el código desde cero, si algún resultado sale raro se puede
ir directamente al código a ver qué está pasando.

\section{Metodología}

El análisis sigue el mismo esquema para los cuatro generadores. Cada uno se configura
con LFSR cuyos polinomios de retroalimentación son primitivos sobre $\mathbb{F}_2$ y
con estados iniciales distintos de cero. Con esa configuración se genera una secuencia
de bits lo bastante larga para que los tests estadísticos sean fiables. Después se
mide la complejidad lineal con Berlekamp-Massey (esto da una idea directa de cuántos
bits necesitaría un atacante para reconstruir el generador) y se pasan los ocho tests
NIST seleccionados. Al final comparo los resultados entre los cuatro generadores para
ver si las diferencias teóricas se reflejan en la práctica.

En concreto, el proceso para cada generador tiene tres fases:

\begin{enumerate}
  \item \textbf{Generación de la secuencia}: se configura el generador con LFSR de
        grados coprimos entre sí (para maximizar el periodo del sistema combinado) y
        se genera una secuencia de 5000 bits. Esta longitud es un compromiso: lo
        bastante larga para que los tests NIST tengan potencia estadística suficiente,
        pero no tanto como para que el cálculo de la complejidad lineal con
        Berlekamp-Massey (que es $O(N^2)$) sea demasiado lento.
  \item \textbf{Análisis de complejidad lineal}: se aplica Berlekamp-Massey a la
        secuencia completa y se obtiene tanto la complejidad lineal final como el
        perfil de complejidad lineal, que muestra cómo evoluciona esta medida a lo
        largo de la secuencia. Una secuencia criptográficamente fuerte debería tener
        un perfil que crece de forma sostenida cerca de la recta $N/2$.
  \item \textbf{Tests estadísticos NIST}: se pasan los ocho tests seleccionados del
        estándar NIST SP 800-22 con un nivel de significación $\alpha = 0{,}01$. Cada
        test calcula un p-valor: si es mayor que $\alpha$, la secuencia supera el test;
        si es menor, se considera que la secuencia tiene alguna debilidad estadística
        detectable por ese test concreto.
\end{enumerate}

Los parámetros concretos (grados de los LFSR, longitud de las secuencias, nivel de
significación de los tests) se detallan en el capítulo~\ref{cap:analisis}. La idea es
mantener condiciones comparables entre generadores para que las diferencias en los
resultados se deban al diseño de cada construcción y no a la configuración experimental.

Los resultados se presentan de dos formas. Para la complejidad lineal se muestra tanto
el valor final como el perfil completo, que permite ver a simple vista si la complejidad
crece de forma regular (como esperaríamos de una secuencia aleatoria) o si se estanca
en valores bajos (lo que indicaría que el generador es vulnerable a Berlekamp-Massey).
Para los tests NIST se genera una tabla resumen con el p-valor de cada generador en
cada test y la indicación de si lo supera o no. Esta presentación uniforme es la que
permite comparar los generadores de forma directa y ver en qué tests concretos falla
cada uno.

\section{Herramientas}

El código está escrito en Python~3. En este trabajo importa más poder leer el código
y verificar que es correcto que exprimir el rendimiento, y las secuencias que genero
no son tan largas como para que el tiempo de ejecución sea un problema.

Como dependencias externas uso \texttt{scipy} para las funciones de distribución que
necesitan los tests estadísticos (en particular la función gamma incompleta y la
distribución chi-cuadrado), \texttt{numpy} para operaciones sobre arrays de bits y
\texttt{matplotlib} para las gráficas de resultados. Los tests unitarios se ejecutan
con \texttt{pytest}, que uso diariamente en el trabajo y permite escribir tests
de forma directa. Para cada módulo del proyecto hay tests unitarios que comprueban
los casos básicos: que un LFSR con polinomio primitivo genera una secuencia de
periodo máximo, que Berlekamp-Massey recupera el polinomio correcto a partir de la
secuencia, que cada generador produce la salida esperada para un estado inicial
conocido, y que cada test NIST da resultados coherentes con secuencias de referencia.

El código fuente está disponible en un repositorio público de GitHub. Para el control
de versiones uso Git con dos ramas de desarrollo separadas: \texttt{feat/code} para
la implementación y \texttt{feat/docs} para la documentación. Esta separación permite
trabajar en el código y en la memoria de forma independiente sin que los cambios de
una parte interfieran con la otra. La estructura del repositorio separa el código
fuente (\texttt{src/}), los tests (\texttt{tests/}), los resultados
(\texttt{results/}) y la documentación (\texttt{docs/}).

La memoria está escrita en \LaTeX{} utilizando la plantilla oficial de la ETSIIT de la
Universidad de Granada. La elección de \LaTeX{} frente a un procesador de textos
convencional viene motivada por la cantidad de notación matemática que tiene el
trabajo: fórmulas en $\mathbb{F}_2$, polinomios, expresiones de complejidad lineal,
tablas de evolución de registros y diagramas. \LaTeX{} maneja todo esto de forma
nativa y produce un resultado tipográfico bastante mejor que el de alternativas como
Word o Google Docs, especialmente cuando hay fórmulas y tablas en casi todas las
páginas.

\section{Alcance}

Este trabajo se centra en la complejidad lineal y las propiedades estadísticas de
cuatro generadores basados en LFSR. No entro en ataques algebraicos ni en los de
\textit{time-memory tradeoff}, que darían para otro trabajo aparte. Tampoco analizo
la seguridad en un escenario de cifrado completo con gestión de claves, ni hay
implementación hardware: lo que me interesa es la calidad criptográfica de las
secuencias, no el rendimiento en un circuito. Tampoco se incluyen construcciones más
recientes como Trivium o Grain, que aunque están basadas en registros de desplazamiento
son diseños considerablemente más complejos y su análisis requeriría herramientas
distintas de las que se desarrollan aquí. El interés de centrarse en los cuatro
generadores clásicos es que son lo bastante sencillos como para entender completamente
su funcionamiento y sus debilidades, y a la vez lo bastante variados como para que la
comparación entre ellos resulte útil.

De los quince tests que define el estándar NIST SP 800-22, implemento ocho. Los
seleccionados son: test de frecuencia (monobit), test de frecuencia dentro de bloques,
test de rachas, test de rachas dentro de bloques, test de rango de matriz binaria,
test espectral (DFT), test de complejidad lineal y test de entropía aproximada. Los
siete restantes (Maurer's universal, cumulative sums, random excursions y sus
variantes, serial y overlapping template matching) se han dejado fuera por distintas
razones: algunos requieren secuencias considerablemente más largas de las que genero,
y otros evalúan propiedades que en el contexto de generadores basados en LFSR resultan
redundantes con los ocho que sí se implementan.

\section{Objetivos}

Los objetivos de este trabajo son:

\begin{itemize}
  \item Implementar desde cero un LFSR genérico y el algoritmo de Berlekamp-Massey
        en Python.
  \item Implementar los cuatro generadores estudiados: Geffe, Shrinking, Self-Shrinking
        y Alternating Step Generator (ASG).
  \item Analizar la complejidad lineal de las secuencias generadas y comprobar si
        coincide con las cotas teóricas.
  \item Pasar ocho tests estadísticos del NIST SP 800-22 a las secuencias de cada
        generador y comparar los resultados entre ellos.
  \item Sacar conclusiones sobre qué generadores ofrecen mejores garantías de seguridad
        y justificar por qué a partir de los resultados obtenidos.
\end{itemize}

\section{Estructura de la memoria}

La memoria se organiza en cinco capítulos. Tras esta introducción, el
capítulo~\ref{cap:fundamentos} presenta los fundamentos teóricos necesarios: la
aritmética en el cuerpo $\mathbb{F}_2$, la definición y propiedades de los LFSR,
las secuencias de longitud máxima, la complejidad lineal y el algoritmo de
Berlekamp-Massey. Se incluyen también los criterios que debe cumplir una buena
secuencia cifrante y un repaso de los principales tipos de ataques a generadores
basados en LFSR.

El capítulo~\ref{cap:generadores} describe los cuatro generadores estudiados: cómo
funciona cada uno, qué propiedades teóricas de seguridad tiene y qué decisiones de
implementación se han tomado. El capítulo~\ref{cap:analisis} presenta los
experimentos: la configuración utilizada, los resultados de complejidad lineal, los
tests NIST y la comparación entre generadores. El último capítulo recoge las
conclusiones del trabajo y las líneas que quedarían abiertas para futuros trabajos.
